%------------------------------------------------------------------------------
% Chinese CV
%------------------------------------------------------------------------------

\documentclass[letterpaper,11pt]{article}
\usepackage{iftex}
\ifXeTeX
  \usepackage{fontspec}
  \setmainfont{Times New Roman}
  \usepackage{xeCJK}
  \IfFontExistsTF{Songti SC}
    {\setCJKmainfont[BoldFont={Songti SC}, BoldFeatures={FakeBold=2}, ItalicFont={Kaiti SC}]{Songti SC}}
    {\IfFontExistsTF{Noto Serif CJK SC}
      {\setCJKmainfont{Noto Serif CJK SC}}
      {\IfFontExistsTF{Hiragino Sans GB}
        {\setCJKmainfont[BoldFont={Hiragino Sans GB W6}]{Hiragino Sans GB}}
        {\setCJKmainfont{WenQuanYi Zen Hei}}}}
\else
  \usepackage[utf8]{inputenc}
  \usepackage[T1]{fontenc}
\fi

\usepackage[empty]{fullpage}
\usepackage{titlesec}
\usepackage[usenames,dvipsnames]{color}
\usepackage{enumitem}
\usepackage[hidelinks]{hyperref}
\usepackage{tabularx}
\usepackage{graphicx}
\usepackage{tikz}
\usetikzlibrary{arrows.meta,calc,positioning}
\usepackage{xcolor}

\addtolength{\oddsidemargin}{-1cm}
\addtolength{\evensidemargin}{-1cm}
\addtolength{\textwidth}{2cm}
\addtolength{\topmargin}{-1cm}
\addtolength{\textheight}{2cm}

\urlstyle{same}
\raggedbottom
\raggedright
\setlength{\tabcolsep}{0cm}

\titleformat{\section}{
  \vspace{-4pt}\bfseries\raggedright\large
}{}{0em}{}[\color{black}\titlerule \vspace{-5pt}]

\newcommand{\CVItem}[1]{\item\small #1 \vspace{-2pt}}
\newcommand{\CVSubheading}[4]{
  \vspace{-2pt}\item
    \begin{tabular*}{0.97\textwidth}[t]{l@{\extracolsep{\fill}}r}
      \textbf{#1} & #2 \\
      \small#3 & \small #4 \\
    \end{tabular*}\vspace{-7pt}
}
\newcommand{\ContactLink}[2]{\href{#1}{\underline{\strut #2}}}
\newcommand{\DOILink}[1]{\href{https://doi.org/#1}{doi:#1}}
\newcommand{\ArxivLink}[1]{\href{https://arxiv.org/abs/#1}{arXiv:#1}}
\newcommand{\ProjectLink}[2]{\href{#1}{#2}}
\newcommand{\CVSubHeadingListStart}{\begin{itemize}[leftmargin=0.5cm, label={}]}
\newcommand{\CVSubHeadingListEnd}{\end{itemize}}
\newcommand{\CVItemListStart}{\begin{itemize}[leftmargin=0.55cm]}
\newcommand{\CVItemListEnd}{\end{itemize}\vspace{-5pt}}
\newcommand{\ResearchCapabilityBox}[5]{
  \node[
    draw=StoryBlue,
    fill=StoryLight,
    rounded corners=2pt,
    text width=4.3cm,
    minimum height=1.9cm,
    align=left,
    inner sep=5pt,
    anchor=north
  ] (#1) at (#2,#3) {
    \centering\color{StoryBlue}\textbf{#4}\par
    \vspace{2pt}
    \raggedright\scriptsize #5
  };
}
\newcommand{\ResearchOutcomeBox}[5]{
  \node[
    draw=StoryGreen,
    fill=StoryGreen!6,
    rounded corners=2pt,
    text width=5.2cm,
    minimum height=1.18cm,
    align=left,
    inner sep=5pt,
    anchor=north
  ] (#1) at (#2,#3) {
    \centering\color{StoryGreen}\textbf{#4}\par
    \vspace{1pt}
    \raggedright\scriptsize #5
  };
}
\newcommand{\ResearchStoryDiagram}{
\begin{center}
\resizebox{\linewidth}{!}{%
\begin{tikzpicture}[font=\scriptsize]
  \definecolor{StoryBlue}{HTML}{0B4AA2}
  \definecolor{StoryLight}{HTML}{EAF1FB}
  \definecolor{StoryGreen}{HTML}{008C73}

  \node[
    draw=StoryBlue,
    fill=StoryLight,
    rounded corners=2pt,
    text width=8.7cm,
    minimum height=0.9cm,
    align=center,
    inner sep=5pt
  ] (core) at (6,0) {\color{StoryBlue}\textbf{面向新型 AI 芯片的}\\\textbf{大模型智能体基础设施}};

  \ResearchCapabilityBox{op}{0.8}{-1.35}{算子功能实现}
    {智能体为大模型负载生成可运行算子，面向新芯片后端落地。}
  \ResearchCapabilityBox{opt}{6}{-1.35}{架构感知优化}
    {智能体结合芯片架构、编译器约束与运行时反馈进行性能优化。}
  \ResearchCapabilityBox{tool}{11.2}{-1.35}{工具链自动开发}
    {智能体开发和修复编译器、运行时与工具链产物。}

  \ResearchOutcomeBox{serving}{3.25}{-3.8}{模型服务效率}
    {更好的算子和运行时路径让新芯片更高效服务大模型。}
  \ResearchOutcomeBox{capability}{8.75}{-3.8}{模型能力增强}
    {开发 trajectory 转化为 coding 与系统智能体训练信号。}

  \draw[-{Latex[length=2mm]}, thick, StoryBlue] (core.south west) -- (op.north);
  \draw[-{Latex[length=2mm]}, thick, StoryBlue] (core.south) -- (opt.north);
  \draw[-{Latex[length=2mm]}, thick, StoryBlue] (core.south east) -- (tool.north);
  \draw[-{Latex[length=2mm]}, thick, StoryGreen] (op.south) -- (serving.north west);
  \draw[-{Latex[length=2mm]}, thick, StoryGreen] (opt.south west) -- (serving.north);
  \draw[-{Latex[length=2mm]}, thick, StoryGreen] (opt.south east) -- (capability.north);
  \draw[-{Latex[length=2mm]}, thick, StoryGreen] (tool.south) -- (capability.north east);

  \node[
    draw=StoryBlue!40,
    fill=StoryBlue!4,
    rounded corners=2pt,
    text width=14.5cm,
    align=left,
    inner sep=5pt,
    anchor=north west
  ] at (-1.35,-5.28) {
    \centering\color{StoryBlue}\textbf{来自既有系统与论文的技术基础}\par
    \vspace{2pt}
    \raggedright\scriptsize
    \textbf{编译器 / 形式化搜索：} \ProjectLink{https://doi.org/10.1145/3779212.3790162}{ISAMORE}、\ProjectLink{https://arxiv.org/abs/2604.17364}{EggMind}\quad
    \textbf{IR / 协同设计接口：} \ProjectLink{https://doi.org/10.1145/3508352.3549370}{Hector}、\ProjectLink{https://doi.org/10.1109/ICCAD66269.2025.11240817}{APS}、\ProjectLink{https://arxiv.org/abs/2511.22267}{Aquas}、\ProjectLink{https://github.com/uv-xiao/IntelliC}{IntelliC}、Spine\par
    \textbf{硬件 / 运行时系统：} \ProjectLink{https://doi.org/10.1145/3626202.3637561}{Cement}、\ProjectLink{https://doi.org/10.1109/ICCAD66269.2025.11240669}{Clay}、\ProjectLink{https://arxiv.org/abs/2511.15323}{SkyEgg}、\ProjectLink{https://doi.org/10.1109/DAC63849.2025.11132875}{Cayman}、\ProjectLink{https://github.com/hw-native-sys/simpler}{PTO Runtime}、\ProjectLink{https://arxiv.org/abs/2604.17353}{Hive}
  };
\end{tikzpicture}
}
\end{center}
}

\begin{document}

\begin{center}
    \textbf{\Huge {肖有为}} \\ \vspace{1pt}
    \small{Youwei Xiao} \\
    \small{+86 185 1920 4005} \\
    \small{\ContactLink{mailto:shallwe@pku.edu.cn}{shallwe@pku.edu.cn} \enspace\textbar\enspace
    \ContactLink{https://uv-xiao.github.io}{uv-xiao.github.io} \enspace\textbar\enspace
    \ContactLink{https://github.com/uv-xiao}{GitHub}}
\end{center}

\section{教育经历}
  \CVSubHeadingListStart
    \CVSubheading
      {北京大学集成电路学院，集成电路科学与工程博士研究生}{2022年9月 -- 至今}
      {导师：梁云教授}{中国北京}
      \CVItemListStart
        \CVItem{研究方向：面向新型 AI 芯片的大模型智能体基础设施，以编译器、运行时与软硬件协同设计为技术基础。}
      \CVItemListEnd
    \CVSubheading
      {北京大学信息科学技术学院，计算机科学与技术理学学士}{2018年9月 -- 2022年7月}
      {}{中国北京}
  \CVSubHeadingListEnd

\section{工作经历}
  \CVSubHeadingListStart
    \CVSubheading
      {大语言模型算法实习生}{2026年 -- 至今}
      {ByteDance Seed}{中国北京}
      \CVItemListStart
        \CVItem{构建面向新型 AI 芯片上大模型系统的智能体基础设施，覆盖算子功能实现、芯片架构感知性能优化与工具链自动开发。}
        \CVItem{设计面向新芯片的大模型服务优化智能体流程，将算子研发、编译器/运行时约束与硬件感知优化连接起来。}
        \CVItem{沉淀真实系统任务中的智能体开发数据与 trajectory，将基础设施研发过程转化为增强 coding 与系统智能体能力的训练信号。}
      \CVItemListEnd
  \CVSubHeadingListEnd

\section{研究方向}

我当前的研究聚焦面向新型 AI 芯片上大模型系统的智能体基础设施，研究智能体如何完成算子功能实现、芯片架构感知性能优化与工具链自动开发。这一方向从两个维度驱动大模型系统：一方面提升新芯片上的模型服务效率，另一方面沉淀开发数据与 trajectory，增强后续 coding 与系统智能体能力。

\ResearchStoryDiagram

\smallskip
\noindent \textbf{大模型智能体基础设施。} 我当前的工作围绕新芯片上大模型基础设施的真实生命周期构建智能体：算子功能实现、架构感知性能优化与工具链自动开发。核心系统问题是：智能体流程如何连接软件意图、编译器/运行时约束、芯片特定优化机会，以及可复用的开发 trajectory。

\smallskip
\noindent \textbf{编译器与运行时底座。} 既有系统工作提供了智能体可操作的系统界面。\textit{IntelliC} (\ProjectLink{https://github.com/uv-xiao/IntelliC}{GitHub}) 研究面向人与智能体协作的可检查编译表示。\textit{PTO Runtime} (\ProjectLink{https://github.com/hw-native-sys/simpler}{GitHub}) 面向 Ascend 芯片和灵衢 SuperPod 执行编译后的任务图。\textit{Hive} (\ProjectLink{https://arxiv.org/abs/2604.17353}{preprint}) 将多智能体大模型负载暴露给推理基础设施，\textit{Spine} 研究跨设计意图、架构、编译器、硬件与执行证据的有界智能体协同设计。

\smallskip
\noindent \textbf{软硬件协同基础。} 我此前在编译器、e-graph、硬件综合与架构定制上的论文构成了当前智能体方向的技术基础。\textit{ISAMORE} (\ProjectLink{https://doi.org/10.1145/3779212.3790162}{DOI}) 和 \textit{EggMind} (\ProjectLink{https://arxiv.org/abs/2604.17364}{preprint}) 展示了形式化搜索与大模型引导搜索如何组织优化过程。\textit{Hector} (\ProjectLink{https://doi.org/10.1145/3508352.3549370}{DOI}, \ProjectLink{https://github.com/pku-liang/Hector}{GitHub})、\textit{Cement} (\ProjectLink{https://doi.org/10.1145/3626202.3637561}{DOI})、\textit{Clay} (\ProjectLink{https://doi.org/10.1109/ICCAD66269.2025.11240669}{DOI})、\textit{Cayman} (\ProjectLink{https://doi.org/10.1109/DAC63849.2025.11132875}{DOI})、\textit{APS} (\ProjectLink{https://doi.org/10.1109/ICCAD66269.2025.11240817}{DOI}, \ProjectLink{https://github.com/pku-liang/APS}{GitHub})、\textit{Aquas} (\ProjectLink{https://arxiv.org/abs/2511.22267}{preprint}) 与 \textit{SkyEgg} (\ProjectLink{https://arxiv.org/abs/2511.15323}{preprint}) 提供了 IR、综合、运行时与协同设计方法，使智能体决策落在真实系统产物上。

\smallskip
\noindent 这条主线使用既有编译器与硬件系统作为基础，迈向下一步：让智能体帮助新芯片更好服务大模型，同时把生成的开发数据与 trajectory 反馈为更强 coding 与系统模型的训练材料。

\section{论文发表}
\CVItemListStart
\CVItem{\textbf{Youwei Xiao}, Chenyun Yin, Yitian Sun, and Yun Liang. 2026. Finding Reusable Instructions via E-Graph Anti-Unification. ASPLOS '26. \DOILink{10.1145/3779212.3790162}. \textbf{\textcolor{red}{[最佳论文奖，5/1048]}}}
\CVItem{\textbf{Youwei Xiao}, Yuyang Zou, Yansong Xu, Yuhao Luo, Yitian Sun, Chenyun Yin, Ruifan Xu, Renze Chen, and Yun Liang. 2025. APS: Open-Source Hardware-Software Co-Design Framework for Agile Processor Specialization. ICCAD '25. \DOILink{10.1109/ICCAD66269.2025.11240817}. [邀请论文]}
\CVItem{Weijie Peng*, \textbf{Youwei Xiao*}, Yuyang Zou, Zizhang Luo, and Yun Liang. 2025. Clay: High-level ASIP Framework for Flexible Microarchitecture-Aware Instruction Customization. ICCAD '25. \DOILink{10.1109/ICCAD66269.2025.11240669}. (*共同一作)}
\CVItem{\textbf{Youwei Xiao}, Fan Cui, Zizhang Luo, Weijie Peng, and Yun Liang. 2025. Cayman: Custom Accelerator Generation with Control Flow and Data Access Optimization. DAC '25. \DOILink{10.1109/DAC63849.2025.11132875}.}
\CVItem{\textbf{Youwei Xiao}, Zizhang Luo, and Yun Liang. 2025. cmt2: Rule-Based Hardware Description in Rust with Temporal Semantics. LATTE '25. \ProjectLink{https://capra.cs.cornell.edu/latte25/paper/1.pdf}{paper}.}
\CVItem{Fan Cui, \textbf{Youwei Xiao}, Kexing Zhou, and Yun Liang. 2025. An Empirical Comparison of LLM-based Hardware Design and High-level Synthesis. FPGA '25. \DOILink{10.1145/3706628.3708861}.}
\CVItem{Fan Cui, Chenyang Yin, Kexing Zhou, \textbf{Youwei Xiao}, Guangyu Sun, Qiang Xu, Qipeng Guo, Yun Liang, Xingcheng Zhang, Demin Song, and Dahua Lin. 2024. OriGen: Enhancing RTL Code Generation with Code-to-Code Augmentation and Self-Reflection. ICCAD '24. \DOILink{10.1145/3676536.3676830}.}
\CVItem{\textbf{Youwei Xiao}, Zizhang Luo, Kexing Zhou, and Yun Liang. 2024. Cement: Streamlining FPGA Hardware Design with Cycle-Deterministic eHDL and Synthesis. FPGA '24. \DOILink{10.1145/3626202.3637561}.}
\CVItem{Ruifan Xu, \textbf{Youwei Xiao}, Jin Luo, and Yun Liang. 2022. HECTOR: A Multi-Level Intermediate Representation for Hardware Synthesis Methodologies. ICCAD '22. \DOILink{10.1145/3508352.3549370}.}
\CVItemListEnd

\section{预印本}
\CVItemListStart
\CVItem{Chenyun Yin*, \textbf{Youwei Xiao*}, Yuze Luo, Yuyang Zou, and Yun Liang. LLM-Guided Strategy Synthesis for Scalable Equality Saturation. \ArxivLink{2604.17364}. (*共同一作)}
\CVItem{Zizhang Luo, Yuhao Luo, \textbf{Youwei Xiao}, Yansong Xu, Runlin Guo, and Yun Liang. Hive: A Multi-Agent Infrastructure for Algorithm- and Task-Level Scaling. \ArxivLink{2604.17353}.}
\CVItem{Yuyang Zou, \textbf{Youwei Xiao}, Yansong Xu, Chenyun Yin, Yuhao Luo, Yitian Sun, Ruifan Xu, Renze Chen, and Yun Liang. Aquas: Enhancing Domain Specialization through Holistic Hardware-Software Co-Optimization based on MLIR. \ArxivLink{2511.22267}.}
\CVItem{\textbf{Youwei Xiao}, Yuyang Zou, and Yun Liang. SkyEgg: Joint Implementation Selection and Scheduling for Hardware Synthesis using E-graphs. \ArxivLink{2511.15323}.}
\CVItem{\textbf{Youwei Xiao}, Zizhang Luo, Weijie Peng, Yuyang Zou, and Yun Liang. Cement2: Temporal Hardware Transactions for High-Level and Efficient FPGA Programming. \ArxivLink{2511.15073}.}
\CVItemListEnd

\section{荣誉奖励}
  \CVSubHeadingListStart
    \CVSubheading
      {学术之芯}{2026}
      {北京大学集成电路学院2026年度荣誉，授予8人}{}
    \CVSubheading
      {博士研究生校长奖学金}{2026-2027}
      {北京大学；集成电路学院授予10人}{}
    \CVSubheading
      {ASPLOS 2026 最佳论文奖}{2026年3月}
      {Finding Reusable Instructions via E-Graph Anti-Unification}{5/1048 submissions}
    \CVSubheading
      {EDAthon 2020 第二名}{2020年8月}
      {电子设计自动化编程竞赛}{IEEE CEDA Hong Kong Chapter}
    \CVSubheading
      {深圳证券交易所奖学金}{2019-2020}
      {北京大学信息科学技术学院17人获奖}{}
    \CVSubheading
      {三好学生}{2018-2019、2019-2020}
      {综合表现优秀，授予前10\%学生}{}
    \CVSubheading
      {杨芙清-王阳元院士奖学金}{2018-2019}
      {北京大学信息科学技术学院仅5人获奖}{}
    \CVSubheading
      {NOI 2017 二等奖}{2017年7月}
      {全国青少年信息学奥林匹克竞赛}{中国计算机学会}
    \CVSubheading
      {APIO 2017 二等奖}{2017年5月}
      {亚太地区信息学奥林匹克中国赛区}{中国计算机学会}
  \CVSubHeadingListEnd

\section{教程报告}
  \CVSubHeadingListStart
    \CVSubheading
      {APS: An MLIR-Based Hardware-Software Co-design Framework}{2026年1月}
      {ASP-DAC 2026}{中国香港}
    \CVSubheading
      {Agile Hardware Specialization: A toolbox for Agile Chip Front-end Design}{2025年5月}
      {ISEDA 2025}{中国香港}
    \CVSubheading
      {Agile Hardware Specialization (AHS)}{2025年3月}
      {ASPLOS 2025}{荷兰鹿特丹}
    \CVSubheading
      {Agile Hardware Specialization: A toolbox for Agile Chip Front-end Design}{2025年3月}
      {DATE 2025}{法国里昂}
    \CVSubheading
      {AHS: An EDA toolbox for Agile Chip Front-end Design}{2025年1月}
      {ASP-DAC 2025}{日本东京}
  \CVSubHeadingListEnd

\section{教学经历}
  \CVSubHeadingListStart
    \CVSubheading
      {高层次芯片设计课程助教}{2022年春季、2022年秋季}
      {北京大学}{}
    \CVSubheading
      {计算机系统导论课程助教}{2020年秋季}
      {北京大学}{}
  \CVSubHeadingListEnd

\end{document}
